% Candidate Figure 1 for the staged-v7 solver. Transparent line art only.
\begin{figure}[H]
\centering
\begin{tikzpicture}[
  every node/.style={font=\fontsize{8.2}{9.6}\selectfont},
  box/.style={rectangle,draw=black,line width=0.45pt,rounded corners=0pt,
    align=center,minimum height=7.2mm,inner xsep=2.5pt,inner ysep=1.5pt},
  wide/.style={box,text width=9.1cm},
  view/.style={box,text width=3.15cm,minimum height=10.8mm},
  review/.style={box,text width=3.15cm,minimum height=9.8mm},
  phase/.style={draw=black,dashed,line width=0.4pt,rounded corners=0pt,
    inner sep=4.5mm},
  arrow/.style={-{Latex[length=2.4mm,width=1.55mm]},line width=0.52pt,
    line cap=rect,line join=miter,shorten >=1pt},
  stem/.style={line width=0.52pt,line cap=rect,line join=miter},
  label/.style={font=\fontsize{8.0}{9.2}\selectfont,inner sep=1pt}
]
\node[wide] (input) at (0,0)
  {读取算例、日期地域参数、有限车队与共同可行初始解};

\node[view] (s1f) at (-4.15,-1.65)
  {HGS-F\\燃油成本视角\\5000次迭代};
\node[view] (s1e) at (0,-1.65)
  {HGS-E\\简化电动视角\\5000次迭代};
\node[view] (s1m) at (4.15,-1.65)
  {HGS-M\\机制感知视角\\5000次迭代};

\node[review] (r1f) at (-4.15,-3.05)
  {完整模型检查\\保留本视角精英};
\node[review] (r1e) at (0,-3.05)
  {完整模型检查\\保留本视角精英};
\node[review] (r1m) at (4.15,-3.05)
  {完整模型检查\\保留本视角精英};

\coordinate (mergeone) at (0,-3.90);
\node[wide] (mipone) at (0,-4.55)
  {汇集可行路线并执行第一次限时MIP组合（30 s）};
\node[wide] (protected) at (0,-5.55)
  {在三视角父方案与可用组合方案中择优，形成阶段1保护解};

\node[view] (s2f) at (-4.15,-7.10)
  {HGS-F\\仅由本视角精英热启动\\继续20000次迭代};
\node[view] (s2e) at (0,-7.10)
  {HGS-E\\仅由本视角精英热启动\\继续20000次迭代};
\node[view] (s2m) at (4.15,-7.10)
  {HGS-M\\仅由本视角精英热启动\\继续20000次迭代};

\node[review] (r2f) at (-4.15,-8.62)
  {完整模型检查\\合并两阶段档案};
\node[review] (r2e) at (0,-8.62)
  {完整模型检查\\合并两阶段档案};
\node[review] (r2m) at (4.15,-8.62)
  {完整模型检查\\合并两阶段档案};

\coordinate (mergetwo) at (0,-9.48);
\node[wide] (miptwo) at (0,-10.13)
  {扩展可行路线池并执行第二次限时MIP组合（30 s）};
\node[wide,minimum height=10.2mm] (judge) at (0,-11.25)
  {完整模型复算与独立检查\\在阶段1保护解、两阶段三视角最优解和可用组合方案中择优};
\node[wide] (output) at (0,-12.42)
  {输出经完整模型核验的最好可行解};

\coordinate (splitone) at (0,-0.76);
\draw[stem] (input.south) -- (splitone);
\draw[arrow] (splitone) -| (s1f.north);
\draw[arrow] (splitone) -- (s1e.north);
\draw[arrow] (splitone) -| (s1m.north);
\draw[arrow] (s1f) -- (r1f);
\draw[arrow] (s1e) -- (r1e);
\draw[arrow] (s1m) -- (r1m);
\draw[stem] (r1f.south) |- (mergeone);
\draw[stem] (r1e.south) -- (mergeone);
\draw[stem] (r1m.south) |- (mergeone);
\draw[arrow] (mergeone) -- (mipone);
\draw[arrow] (mipone) -- (protected);

\coordinate (splittwo) at (0,-6.30);
\draw[stem] (protected.south) -- (splittwo);
\draw[arrow] (splittwo) -| (s2f.north);
\draw[arrow] (splittwo) -- (s2e.north);
\draw[arrow] (splittwo) -| (s2m.north);
\draw[arrow] (s2f) -- (r2f);
\draw[arrow] (s2e) -- (r2e);
\draw[arrow] (s2m) -- (r2m);
\draw[stem] (r2f.south) |- (mergetwo);
\draw[stem] (r2e.south) -- (mergetwo);
\draw[stem] (r2m.south) |- (mergetwo);
\draw[arrow] (mergetwo) -- (miptwo);
\draw[arrow] (miptwo) -- (judge);
\draw[arrow] (judge) -- (output);

\begin{scope}[on background layer]
  \node[phase,fit=(s1f)(s1m)(r1f)(r1m)] (phaseone) {};
  \node[phase,fit=(s2f)(s2m)(r2f)(r2m)] (phasetwo) {};
\end{scope}
\node[label,above=0.5mm of phaseone.north]
  {阶段1：基准搜索与保护};
\node[label,above=0.5mm of phasetwo.north]
  {阶段2：同视角加深搜索};
\end{tikzpicture}
\caption{两阶段多视角混合遗传搜索流程}
\label{fig:algorithm-flow}
\end{figure}
